\documentclass[tikz,border=4pt]{standalone}
\usepackage{tikz}
\usepackage{polymer-paper-preamble}

\tikzset{
  >={Stealth[length=3.2pt,width=2.4pt]},
  panel title/.style={font=\sffamily\bfseries\small, text=black, anchor=west},
  label/.style={font=\sffamily\footnotesize, text=cGray},
  compact label/.style={font=\sffamily\fontsize{6.5}{7.5}\selectfont, text=cGray},
  axis/.style={draw=cGray, line width=0.55pt, -{Stealth}},
  guide/.style={draw=cGray!70, line width=0.35pt},
  carrier/.style={circle, draw=cBlue, fill=white, line width=0.55pt,
                  minimum size=2.5mm, inner sep=0pt},
  site/.style={draw=cRed, line width=0.55pt}
}

\begin{document}
\begin{tikzpicture}[x=1cm,y=1cm]

% Panel A: applied bias, dispersive repeated trapping, and qualitative decay.
\node[panel title] at (0.00,4.42) {(a) Trapping during conduction};

\path[draw=cGray, fill=cAmberSphere!22, line width=0.65pt, rounded corners=1.5pt]
  (0.18,1.12) rectangle (3.66,3.76);
\path[draw=cGray, fill=cGray, line width=0.55pt]
  (0.18,1.12) rectangle (0.47,3.76);
\path[draw=cGray, fill=cGray, line width=0.55pt]
  (3.37,1.12) rectangle (3.66,3.76);
\node[font=\sffamily\fontsize{6.5}{7.5}\selectfont, text=white, rotate=90]
  at (0.325,2.44) {electrode};
\node[font=\sffamily\fontsize{6.5}{7.5}\selectfont, text=white, rotate=90]
  at (3.515,2.44) {electrode};
\node[compact label, align=center, anchor=north] at (1.92,3.83)
  {disordered\\sulfur--polymer film};

\draw[draw=cGray, line width=0.45pt]
  (0.325,3.76) -- (0.325,4.04) -- (3.515,4.04) -- (3.515,3.76);
\node[label, fill=white, inner sep=1.3pt] at (1.92,4.04) {applied $V$};

% Unassigned trap-site marks and sparse disorder cues.
\foreach \x/\y in {0.78/1.52,0.94/2.75,1.34/2.02,1.61/3.03,
                    2.05/1.62,2.31/2.59,2.68/1.98,3.02/3.02}
  \draw[site] (\x-0.055,\y-0.055)--(\x+0.055,\y+0.055)
              (\x-0.055,\y+0.055)--(\x+0.055,\y-0.055);
\foreach \x/\y/\r in {0.68/2.22/0.045,1.14/3.20/0.055,1.52/1.42/0.045,
                        1.87/2.28/0.055,2.45/3.25/0.045,2.83/1.38/0.055,
                        3.15/2.35/0.045}
  \fill[cAmber] (\x,\y) circle (\r);

% The walk changes direction repeatedly and has no electrode-polarity assignment.
\draw[draw=cBlue, line width=0.95pt, rounded corners=2.5pt]
  (1.08,1.62) -- (1.34,2.02) -- (1.07,2.42) -- (1.61,3.03)
  -- (2.02,2.72) -- (2.31,2.59) -- (2.02,2.18) -- (2.68,1.98)
  -- (2.49,1.52) -- (3.02,3.02);

\node[carrier] at (1.08,1.62) {};
\node[carrier] at (1.34,2.02) {};
\node[carrier] at (2.31,2.59) {};
\node[carrier, fill=cBlue!18] at (3.02,3.02) {};
\node[compact label, align=center, anchor=north] at (1.92,1.03)
  {capture $\leftrightarrow$ release repeats\\filled endpoint: retained};
% Transition and endpoint semantics share one external key.

% Qualitative inset only: no ticks, values, or sampled data markers.
\begin{scope}[shift={(4.05,1.40)}]
  \draw[axis] (0,0)--(1.62,0);
  \draw[axis] (0,0)--(0,1.65);
  \node[compact label] at (0.81,-0.18) {$t$};
  \node[compact label, rotate=90] at (-0.22,0.825) {$I(t)$};
  \draw[draw=cTeal, line width=1.00pt]
    (0.10,1.46) .. controls (0.27,0.87) and (0.72,0.42) .. (1.50,0.20);
  \node[compact label, text=cTeal, align=left, anchor=west]
    at (0.43,1.28) {qualitative\\$I(t)\propto t^{-n}$};
  \node[label, align=center] at (0.81,-0.43)
    {$I(t)\downarrow\;\Longrightarrow\;R\uparrow$};
\end{scope}

\draw[draw=cGray!70, line width=0.45pt, -{Stealth}] (5.82,2.46)--(6.25,2.46);

% Panel B: composition-dependent trap-energy distribution.
\node[panel title] at (6.37,4.42) {(b) Trap-energy landscape};
\node[label, anchor=west] at (7.25,4.03) {sulfur content increases};
\draw[draw=cAmber, line width=0.55pt, -{Stealth}] (10.55,4.03)--(11.36,4.03);

\begin{scope}[shift={(6.58,0.86)}]
  \draw[axis] (0,0)--(4.92,0) node[compact label, below left=0pt]
    {trap energy, $E$};
  \draw[axis] (0,0)--(0,2.77);
  \node[compact label, anchor=south east] at (-0.05,2.82) {$g(E)$};

  % S60 is represented as one narrow, discrete state.
  \path[fill=cBlue!12, draw=none]
    (0.44,0) .. controls (0.68,0.03) and (0.75,0.12) .. (0.86,0.90)
    .. controls (0.93,1.42) and (1.00,1.42) .. (1.07,0.90)
    .. controls (1.18,0.12) and (1.25,0.03) .. (1.50,0) -- cycle;
  \draw[draw=cBlue, line width=1.00pt]
    (0.44,0) .. controls (0.68,0.03) and (0.75,0.12) .. (0.86,0.90)
    .. controls (0.93,1.42) and (1.00,1.42) .. (1.07,0.90)
    .. controls (1.18,0.12) and (1.25,0.03) .. (1.50,0);
  \draw[draw=cBlue!70, line width=0.35pt, dashed] (0.965,0)--(0.965,1.30);
  \node[label, text=cBlue, font=\sffamily\bfseries\footnotesize] at (0.965,2.06) {S60};
  \node[compact label] at (0.965,1.68) {single discrete state};

  % S80 is a broad continuous distribution; width, not depth, carries n.
  \path[fill=cRed!14, draw=none]
    (2.10,0) .. controls (2.32,0.10) and (2.50,0.63) .. (2.82,1.05)
    .. controls (3.18,1.49) and (3.57,1.49) .. (3.91,1.04)
    .. controls (4.22,0.63) and (4.40,0.10) .. (4.63,0) -- cycle;
  \draw[draw=cRed, line width=1.00pt]
    (2.10,0) .. controls (2.32,0.10) and (2.50,0.63) .. (2.82,1.05)
    .. controls (3.18,1.49) and (3.57,1.49) .. (3.91,1.04)
    .. controls (4.22,0.63) and (4.40,0.10) .. (4.63,0);
  \node[label, text=cRed, font=\sffamily\bfseries\footnotesize] at (3.36,1.92) {S80};
  \node[compact label] at (3.36,1.65) {continuous broad};

  % Orthogonal metrics: horizontal breadth and detached vertical magnitude.
  \draw[draw=cTeal, line width=0.65pt, -{Stealth}] (2.75,0.28)--(2.23,0.28) (3.98,0.28)--(4.50,0.28);
  \node[compact label, text=cTeal, fill=white, inner sep=1.4pt] at (3.365,0.28)
    {$n$: breadth};

  \draw[draw=cBrown, line width=0.65pt, {Stealth}-{Stealth}] (4.78,0.18)--(4.78,1.35);
  \node[compact label, text=cBrown, anchor=west] at (4.93,1.24) {$\rho_{60\mathrm{s}}$};
  \node[compact label, text=cBrown, rotate=90] at (4.58,0.765) {magnitude};
\end{scope}

\node[label, align=center] at (8.94,0.18)
  {discrete S60 $\longrightarrow$ continuous broad S80};

\end{tikzpicture}
\end{document}
